\begin{abstract}
Judgmental forecasting requires inferring predictive relationships from sparse
observations and contextual knowledge. We ask whether language models recover
and use these relationships. Across prediction markets and synthetic worlds,
models selectively update forecasts, infer relationships from context, and
revise them as direct evidence accumulates. Larger tested models infer arbitrary
conventions; intervening on one model's internal representation changes its
forecast. Training improves forecasts across domains, but an eight-seed
replication weakens structural-transfer claims. A supervised pilot learns
response-form selection in distribution where the tested RL recipe fails,
without transferring to a new domain or arbitrary labels. Accuracy gains on
compositions of familiar mechanisms exceed a fixed prior, but a gain intervention
does not establish transferred evidence use. Historical-market training lowers
mean error on unseen event families. These findings support context-sensitive
inductive forecasting while distinguishing in-domain learnability, transfer,
and episode-specific evidence use.\footnotemark
\end{abstract}
\footnotetext{%
  \micon{huggingface.png}\;%
  \href{https://huggingface.co/datasets/anonymous-review/forecasting-study-data}%
  {\textbf{Study dataset and artifacts}}\qquad
  \micon{gh_logo.png}\;%
  \href{https://github.com/anonymous-review/forecasting-study-code}%
  {\textbf{Code repository}}\quad
  \textit{Links anonymized for review.}}

\section{Introduction}
\label{sec:introduction}

\setlength{\columnsep}{6pt}
\begin{wrapfigure}{r}{0.44\textwidth}
\vspace{-1.2em}
\centering
\includegraphics[width=\linewidth,trim=0 72pt 0 0,clip]{teaser_revised.pdf}
\vspace{-0.45cm}
\caption{\textbf{Motivation.} Judgmental forecasts combine observations with
contextual knowledge about the process that generated them. Frontier models,
asked to forecast this question, average $P(H)\approx0.68$ (see
App.~\ref{app:carnival-coin-probe}).}
\vspace{0pt}
\label{fig:world-knowledge-evidence}
\end{wrapfigure}

Judgmental forecasting is a demanding form of reasoning central to decisions
about political, social, and economic futures. Unlike settings with abundant
historical data for statistical extrapolation, judgmental forecasting requires
integrating sparse, noisy evidence with contextual knowledge to form calibrated
beliefs about future outcomes
\citep{tetlock2016superforecasting,mellers2014psychological,lawrence2006judgmental}.
This creates an inductive challenge: forecasters must
infer from limited observations the relationships that govern how evidence
impacts outcomes, and use those relationships to anticipate how events may
unfold
\citep{tenenbaum2006theory,griffiths2010probabilistic,tenenbaum2011mind,lake2017building}.
Figure~\ref{fig:world-knowledge-evidence} illustrates this identification problem
in a low-stakes setting. A researcher observes a carnival coin-flipping machine
and, after seeing two heads, must estimate the probability that the next flip
will be heads. A judgmental forecast may incorporate contextual knowledge (e.g.,
whether the game is rigged) to infer the process generating the observed
outcomes.

Language models (LMs) now approach crowd accuracy in judgmental forecasting
\citep{halawi2024approaching,schoenegger2024wisdom}, but accurate predictions need
not imply recovery of the underlying task structure. Models recover latent
structure in some settings \citep{li2023emergent}, yet can predict accurately
using relationships that fail to generalize in others \citep{vafa2025foundation}.
This distinction matters when decision-makers ask how and whether a forecast
should change under alternative evidence or events. Answering such questions
requires predictive relationships that generalize across changes in context.

\noindent
\textbf{Present work.}
We study whether LMs exhibit inductive reasoning in forecasting tasks through a
series of progressively more stringent behavioral tests spanning both real-world
prediction markets and controlled synthetic environments. First, we evaluate
whether LMs respond selectively to new evidence, distinguishing
outcome-relevant information from topically related but orthogonal information
(\hyperref[sec:exp1]{Exp.~1}). Second, we test whether LMs infer predictive
relationships from described context before statistical evidence accumulates. We
construct synthetic forecasting problems in which the numerical observations are
identical across conditions and only the described setting varies, isolating the
contribution of background knowledge to the LM's forecast
(\hyperref[sec:exp2]{Exp.~2}). Finally, we test whether relationships learned
during fine-tuning transfer to new settings. We fine-tune LMs with reinforcement
learning from verifiable rewards (RLVR) on both synthetic tasks and historical
prediction market data, diagnose component failures with a supervised pilot,
and evaluate on held-out structures and domains
(\hyperref[sec:exp3]{Exp.~3}, \hyperref[sec:exp4]{Exp.~4}).

\noindent\textbf{Results.}
All nine systems move $1.9$--$7.5\times$ farther for outcome-relevant than topical
evidence, conditional on updating (\hyperref[sec:exp1]{Exp.~1}). In Coin City,
forecasts track context-selected relationships ($r=.52$--$.70$) and revise them
as direct evidence arrives; larger models infer episode-local conventions, with
a causal intervention linking the induced representation to Qwen3-14B's forecast
(\hyperref[sec:exp2]{Exp.~2}). Training transfers across domains, but its
advantage over an answer-matched shuffled control under joint domain and structure
shift is uncertain at eight seeds (Qwen3-8B:
$\eEightJointEst{}$, 95\% CI $\eEightJointCi{}$). A supervised pilot learns response-form selection in distribution but fails
under domain or label changes. A separate composition benchmark improves
accuracy over a fixed prior, without clear evidence of transferred gain tracking
(\hyperref[sec:exp3]{Exp.~3}). On 318 held-out markets, training lowers mean
error for every tested base and several checkpoints reach crowd performance,
although the Qwen3-32B improvement is uncertain
(\hyperref[sec:exp4]{Exp.~4}).
